%----------------------------------------------------------------------------------------
% Tailored CV — Clinical Bioinformatics PhD, CiiM / MHH (Prof. Cheng-Jian Xu)
%----------------------------------------------------------------------------------------
\documentclass[11pt,a4paper]{article}

\usepackage[utf8]{inputenc}
\usepackage[T1]{fontenc}
\usepackage[margin=2.5cm]{geometry}
\usepackage{hyperref}
\usepackage{microtype}
\usepackage{parskip}
\usepackage{lmodern}

\hypersetup{colorlinks=true, urlcolor=blue!60!black}

\pagestyle{empty}

\begin{document}

\begin{flushright}
Kundai Farai Sachikonye \\
Auf der Lichtung 19, 82178 Puchheim, Germany \\
\href{https://github.com/fullscreen-triangle}{github.com/fullscreen-triangle} \\
\texttt{kundai.sachikonye@bitspark.com} \\
(+49) 1626386344 \\[6pt]
\today
\end{flushright}

\quote{Clinical bioinformatician: multi-omics integration, disease-state modelling,
and predictive stratification --- reproducible, validated, and at population scale.}

%----------------------------------------------------------------------------------------

\begin{document}

\makecvtitle

%----------------------------------------------------------------------------------------
%	PROFILE
%----------------------------------------------------------------------------------------

\section{Profile}
\cvitem{}{
  Computational biologist working on integrated multi-omics analysis and formal
  models of disease state across biological scales. I build reproducible,
  HPC-scale pipelines validated against authoritative reference data (gnomAD,
  UK~Biobank, IEDB, NIST), and low-dimensional disease-vector models for
  multimorbidity, trajectory, and patient stratification. Fully authorised to
  work in Germany without restriction (permanent residence).
}

%----------------------------------------------------------------------------------------
%	EDUCATION
%----------------------------------------------------------------------------------------

\section{Education}

\cventry{2019--2021}{PhD candidate, Computational Lipidomics}{Technische Universität München}{}{}{
  Doctoral research in mass spectrometry-based lipidomics, multi-omics statistical
  modelling, and scientific computing. Departed programme to pursue independent
  computational research.
}

\cventry{2018--2019}{MSc Computational Biology}{Universität Tübingen}{}{}{}

\cventry{2014--2017}{MSc Pharmaceutical Biotechnology}{HAW Hamburg}{}{}{}

\cventry{2011--2014}{BSc Biochemistry and Cell Biology}{Jacobs University Bremen}{}{}{}

%----------------------------------------------------------------------------------------
%	RESEARCH OUTPUT
%----------------------------------------------------------------------------------------

\section{Research Output}

\cvitem{2025}{
  \textbf{Syndrome: Categorical Resolution of Disease Trajectories.}
  \newline
  Disease state encoded as a vector $\mathbf{D}\in[0,1]^8$ over eight cellular
  oscillator classes (protein, enzyme, channel, membrane, ATP, genetic expression,
  calcium, circadian); multimorbidity, trajectory, and stratification become
  geometry on a shared manifold. Intervention target $\arg\max_i w_i D_i$ as a
  sparse $\ell_1$ LP with a priori side-effect bound; $\mathcal{O}(N\log N)$
  trajectory resolution. \textit{Manuscript.}
}

\cvitem{2025}{
  \textbf{On Disease Epistemology in Blind Receiver: Self-Consistency as the
  Unique Basis of Disease Detection.}
  \newline
  Proves template-based diagnosis is epistemologically insufficient; derives
  loop-holonomy self-consistency as the unique viable criterion
  ($H_\ell\neq\mathrm{Id}$ signals disease). Detects distributed, circuit-level
  inconsistency invisible to single-marker panels --- relevant to gradual
  processes such as immune aging. AUC$=1.00$ vs 0.896 for template comparison;
  25/25 validation experiments pass. \textit{Manuscript.}
}

\cvitem{2025}{
  \textbf{On the Thermodynamic Consequences of Categorical Completion on
  Biological Systems (adaptive immunity).}
  \newline
  MHC binding affinity correlates with parent-protein categorical richness
  independently of classical sequence features ($r=0.73$, $p<10^{-12}$,
  $n=2{,}847$ IEDB peptides). Adding the term to NetMHCpan improves AUC
  0.68~$\to$~0.81 ($p<0.001$, DeLong). \textit{Manuscript.}
}

%----------------------------------------------------------------------------------------
%	RESEARCH SOFTWARE
%----------------------------------------------------------------------------------------

\section{Research Software}

\cvitem{gospel}{
  Multi-omics integration pipeline: whole-genome VCF integrated with RNA-seq
  expression at $10^6$ variants/second; Bayesian network inference, Gaussian
  mixture models, variational Bayes. Validated against 1000~Genomes Phase 3,
  gnomAD v3.1.2, UK~Biobank.
  \href{https://github.com/fullscreen-triangle/gospel}{github.com/fullscreen-triangle/gospel}
}

\cvitem{lavoisier}{
  Mass-spectrometry metabolomics framework; 98.9\% feature-extraction accuracy
  against NIST reference libraries; PyTorch models, distributed computing (Ray),
  $>$100\,GB mzML datasets via memory-mapped I/O.
  \href{https://github.com/fullscreen-triangle/lavoisier}{github.com/fullscreen-triangle/lavoisier}
}

\cvitem{syndrome}{
  Implementation of the disease-vector framework: oscillator coherence
  computation, immunopeptidomics (MHC prediction, AUC~0.81 on IEDB), trajectory
  resolution. Python/Rust; browser-deployable tools.
  \href{https://syndrome-seven.vercel.app/tools}{syndrome-seven.vercel.app}
}

\cvitem{helicopter}{
  Multi-modal microscopy image analysis (fluorescence, brightfield, spectral)
  via sequential exclusion.
  \href{https://github.com/fullscreen-triangle/helicopter}{github.com/fullscreen-triangle/helicopter}
}

%----------------------------------------------------------------------------------------
%	RESEARCH EXPERIENCE
%----------------------------------------------------------------------------------------

\section{Research Experience}

\cventry{2019--2021}{Computational Lipidomics --- Bioinformatics \& HPC}{Technische Universität München}{Munich}{}{
  High-throughput mass-spectrometry bioinformatics at HPC scale ($>$100\,GB
  datasets): peak detection, ion database annotation, feature engineering,
  statistical modelling, multi-omics visualisation. Federated learning
  architecture for distributed analysis across clinical sites. Taught at
  Master's level.
}

\cventry{2017--2018}{Glycomics and Proteomics}{Universitätsklinikum Hamburg-Eppendorf}{Hamburg}{}{
  Automated container-based pipeline for standardised LC-MS/MS and MALDI-TOF
  characterisation; high-throughput processing and statistical QC.
}

\cventry{2013}{Cancer Biomarker Discovery}{University Medical Center Groningen}{Groningen}{}{
  UPLC-MS/MS shotgun proteomics protocol optimisation; high-dimensional data
  analysis for cervical cancer biomarker discovery.
}

\cventry{2013}{Pharmacokinetics}{Leibniz Lungenzentrum}{Borstel}{}{
  Quantification of matrix suppression in anti-TB drug efficacy measurement
  (Selected Reaction Monitoring LC-MS).
}

%----------------------------------------------------------------------------------------
%	TECHNICAL SKILLS
%----------------------------------------------------------------------------------------

\section{Technical Skills}

\cvitem{Languages}{Python, R, Rust, TypeScript/JavaScript, Java, C, SQL, Bash}

\cvitem{Multi-omics}{
  Variant calling (VCF), RNA-seq / transcriptome analysis, single-cell RNA-seq,
  mass-spectrometry metabolomics/proteomics (mzML, mzXML, BAM, FASTA),
  cross-omics integration, pathway analysis (KEGG, Reactome)
}

\cvitem{Statistics / ML}{
  Bayesian networks, Gaussian mixture models, variational Bayes, Gaussian
  processes, gradient boosting, graph neural networks, Transformers; PyTorch,
  scikit-learn, NumPy, SciPy, Pandas
}

\cvitem{Infrastructure}{
  Docker, distributed computing (Ray, Dask), HPC / SLURM, cloud (AWS, Azure),
  Git, PostgreSQL, Neo4j, Jupyter
}

\cvitem{Reference data}{
  1000~Genomes, gnomAD, UK~Biobank, IEDB, NIST
}

%----------------------------------------------------------------------------------------
%	LANGUAGES
%----------------------------------------------------------------------------------------

\section{Languages}

\cvitemwithcomment{English}{Native / Fluent (C2)}{}
\cvitemwithcomment{German}{Conversational}{resident in Germany since 2011}

\renewcommand{\listitemsymbol}{-~}

\end{document}
